\section{Analysis and Conclusions}

This chapter summarizes the performance, efficiency, and applicable scenarios of each architecture and identifies the \textit{performance hero} and \textit{efficiency hero} of the project.

Under the current ZU4EV mainline, both the \textit{performance hero} and the \textit{efficiency hero} are won by \texttt{fir\_pipe\_systolic}. It not only achieves the highest throughput, but is also the only custom architecture that closes timing at the target period of \texttt{3.333 ns}. This indicates that for a project of this class, with 261 taps, linear phase, and strict bit-true FIR behavior, a high-quality symmetry-folded systolic pipeline remains a very strong engineering solution.

The research value of \texttt{L=2} and \texttt{L=3} still stands. \texttt{L=2} shows how a true polyphase + symmetry structure can enter the unified verification and implementation chain; \texttt{L=3} proves that the compressed route with a shared \texttt{L3 FFA core} already has presentation-worthy high-throughput potential on ZU4EV. At the same time, the system-level \texttt{PS + PL + Bare-metal Vitis} flow is also closed: both \texttt{fir\_pipe\_systolic} and \texttt{vendor FIR IP} have completed automatic download, automatic UART capture, and on-board case validation, with \texttt{mismatch\_count = 0} on board.

The final project conclusion must be presented honestly under two framings. If we discuss the custom research RTL matrix, the champion is still \texttt{fir\_pipe\_systolic}; if we discuss the final system-shell comparison, \texttt{vendor FIR IP} achieves lower LUT, FF, and \texttt{energy/sample} while keeping almost the same frequency. Therefore, the value of this project is not ``forcing a proof that the custom design wins everywhere,'' but rather that under a unified specification, unified fixed-point model, unified vectors, and unified board-validation chain, we have placed the custom architecture and the industrial baseline onto the same reproducible comparison table.

This round also filled in two analytical dimensions that matter strongly for the final conclusion. The first is stability: both \texttt{fir\_pipe\_systolic} and \texttt{vendor FIR IP} have now completed the latest \texttt{3} repeated windows under the current formal \texttt{8-case} board-validation suite, and all \texttt{8/8} cases passed with \texttt{mismatch\_sum = 0}. This shows that the current automated closure is not just ``a single successful run,'' but already has stable repeat-run capability. The second is methodological transparency: the repository now preserves not only the final numbers, but also the methodological framing, power breakdown, critical-path decomposition, bit-width derivation, and stability statistics as CSV/JSON sources of truth, so that both GitHub Pages and the course report can generate consistent figures and tables directly from those sources.
